\documentclass[11pt]{article}
\usepackage[margin=1in]{geometry}
\usepackage{amsmath,amssymb,amsthm}
\usepackage[T1]{fontenc}
\usepackage{hyperref}
\newcommand{\NE}{\operatorname{NE}}
\newcommand{\Stab}{\operatorname{Stab}_0}
\newcommand{\Reg}{\operatorname{Reg}}
\newcommand{\SwapReg}{\operatorname{SwapReg}}
\newcommand{\pos}[1]{[#1]_+}
\newtheorem{proposition}{Proposition}
\title{Constant-regret literature update and the remaining catalogue gap}
\author{}
\date{Status checked: 17 October 2026}
\begin{document}
\maketitle
\section*{Target and claim}
This is a literature report for \texttt{ECON-C73-63},
\emph{Minimax self-play regret under adversarially valid decentralized
learning}. No new algorithm or full solution is claimed.
\section*{Game, feedback, and regret}
Fix integers $N\geq 2$ and $K\geq 2$. A game has labelled players
$i\in\{1,\ldots,N\}$, action sets $A_i=\{1,\ldots,m_i\}$ with
$1\leq m_i\leq K$, and payoffs $u_i:\prod_j A_j\to[0,1]$.
Payoffs extend multilinearly to $\Sigma=\prod_i\Delta(A_i)$.
At round $t\geq1$, all players choose mixed strategies simultaneously.
After play, player $i$ receives its own exact vector
$v_i^t(a)=u_i(a,x_{-i}^t)$ for every $a\in A_i$.
An uncoupled update uses its own past vectors, its state, its own label and
the public dimensions, without access to any player's payoff table or to
precomputed equilibria. All logarithms are natural. Set
\[
 R_i(T)=\max_{a\in A_i}\sum_{t=1}^T
  \bigl(v_i^t(a)-\langle x_i^t,v_i^t\rangle\bigr),
 \qquad [r]_+=\max\{r,0\}.
\]
The comparator is one fixed action. A horizon-uniform bound concerns
$\sup_T R_i(T)$, rather than $\sup_T|R_i(T)|$ or a comparator that changes
actions each round. This is expected-payoff feedback and mixed-payoff
regret, rather than realized regret after action sampling.

\section*{Reported advances}
Liu--Farina--Ozdaglar introduce ECHO-OFTRL and report deterministic,
uncoupled, horizon-independent individual regret; their dimensional
rate is $O(N^{21}\log^4K)$ as compared in HOOD's introduction
\cite{echo,hood}. Abbadi--Laraki--Mertikopoulos report for HOOD
\[
 R_i(T)\leq45(N+1)^3(16+4\log K+\log^2K),\qquad i,T,
\]
in the same expected-payoff self-play model \cite{hood}.
Soleymani--Piliouras report the stronger dimensional bound
\cite[Theorem 4.10]{morm}
\[
 R_i(T)\leq A(N,K):=96\sqrt N(2+\log K),\qquad i,T,
\]
with constant rate $\eta=1/(32\sqrt N)$. Their safeguarded rule leaves
fixed-game self-play unchanged and has the adversarial guarantee
\[
 R_i(T)\leq A(N,K)+21\sqrt{T(2+\log K)}.
\]
These are attributed preprint theorem claims, not independently
certified proofs or results originating in this contribution.

\section*{Why the existing target needs a separate admissibility argument}
The catalogue fixes $C_0$ and a declared polynomial per-round budget.
For each player, admissibility requires expected positive regret at most
$C_0\sqrt{T\log m_i}$ against every allowed nonanticipating payoff-vector
sequence. Its minimax quantity is
\[
 \mathcal R_{N,m}(T)=\inf_{\mathcal A\in\mathfrak L(C_0)}\sup_G
 \max_i\mathbb E[(R_i^{\mathcal A,G}(T))_+].
\]
An adversarial rate with a dimension-dependent additive constant is not,
by itself, a proof of membership in this fixed class. In particular, the
displayed MORM bound does not identify the minimum prescribed $C_0$ for
which that membership holds. Any use of an exact-arithmetic algorithm
must also match the target's declared computational model. These
requirements cannot be dropped when importing the new literature.

For illustration, using the trivial pathwise bound $R_i(T)\leq T$,
\[
 [R_i(T)]_+\leq
 \min\{T,A+21\sqrt{T(2+\log K)}\}.
\]
With $L_i=\log m_i>0$, $\min\{T,A\}\leq\sqrt{AT}$ implies
\[
 [R_i(T)]_+\leq
 \left(\sqrt{A/L_i}+21\sqrt{(2+\log K)/L_i}\right)
 \sqrt{TL_i}.
\]
This elementary comparison supplies a sufficient, dimension-dependent
budget coefficient for the reported safeguarded bound, without proving
the required computational admissibility or covering a smaller fixed
$C_0$. Players with $m_i=1$ have regret zero. It is not a matching lower
bound or a proof of optimality.

\section*{Existing contributions and relation}
The three existing contributions were read at the pinned source revision
recorded in the manifest. Attempt A1 gives an initial-information lower
bound and a constant-vector specialization. A2 extends a wrapper to a
persistent maximizing action and weak-dominance situations. A3 treats
iterated strict dominance with game-dependent constants and discusses
finite precision. Their stated scopes do not give a dimension-uniform
bound over all games; none is superseded by a claim of a full solution
here. This report adds the general-game preprint context.

\section*{Equilibrium consequence and scope}
For $\mu_T=T^{-1}\sum_{t=1}^T\bigotimes_jx_j^t$, define the nonnegative
CCE gap by taking the positive part of the maximum fixed-action gain.
Linearity gives
\[
 \operatorname{Gap}_{\rm CCE}(\mu_T)
   =T^{-1}\max_i[R_i(T)]_+.
\]
This concerns the time mixture of product distributions. It does not
prove last-iterate Nash convergence, constant realized-action regret,
or the Markov-game target \texttt{ECON-C73-113}. The stochastic-game
payoff target \texttt{ECON-C73-1} is also a different question.

\section*{Preparation and attribution}
The original research map and analysis were generated by OpenAI GPT 6
Astra Pro under the research direction and prompting of GitHub user
\texttt{drmehmetismail}. An OpenAI Codex assistant based on GPT-6 checked
the cited source statements, compared the catalogue scope, and prepared
this LaTeX literature report. Exact model versions and reasoning settings
were not provided. No human mathematical authorship or independent human
proof verification is claimed. Cited researchers receive credit for their
own published or preprint results, rather than authorship of this report.
Only source statements and scope were checked; no literature-wide novelty
certification or complete audit of the cited proofs is claimed.
\begin{thebibliography}{9}
\bibitem{echo}
Mingyang Liu, Gabriele Farina, and Asuman Ozdaglar.
\emph{Constant Individual Regret in General Games}.
arXiv:2608.31166v1, 31 August 2026.
Abstract; comparison in HOOD introduction.
\url{https://arxiv.org/abs/2608.31166v1}.
\bibitem{hood}
Omar Abbadi, Rida Laraki, and Panayotis Mertikopoulos.
\emph{Constant Regret in General Games via Higher-Order Optimism}.
arXiv:2609.04113v1, 3 September 2026.
Introduction, informal main theorem; full-information setup.
\url{https://arxiv.org/html/2609.04113v1}.
\bibitem{morm}
Ashkan Soleymani and Georgios Piliouras.
\emph{Multiplicative Optimism for Constant Regret in Games}.
arXiv:2609.21976v1, 18 September 2026.
Theorem 4.10; Corollary 2.5; Theorem E.1; Proposition E.2; conclusion.
\url{https://arxiv.org/html/2609.21976v1}.
\end{thebibliography}
\end{document}
